\chapter{Resultados Preliminares}
\label{sec:resultados}

Este capítulo apresenta os resultados preliminares obtidos a partir da aplicação inicial da metodologia proposta. O objetivo desta etapa é avaliar a viabilidade do estudo, identificar tendências iniciais e fornecer evidências preliminares que sustentem a continuidade da pesquisa.

\section{Configuração Experimental Inicial}
\label{subsec:config_preliminar}

\textbf{Ideias a abordar:}
\begin{itemize}
    \item Descrição do experimento realizado:
    \begin{itemize}
        \item quantidade de questões utilizadas
        \item origem do dataset (mesmo que parcial)
    \end{itemize}
    \item Modelos avaliados:
    \begin{itemize}
        \item quais LLMs foram utilizados
    \end{itemize}
    \item Estratégias de prompting testadas:
    \begin{itemize}
        \item zero-shot
        \item outras variações (se aplicável)
    \end{itemize}
    \item Escopo da análise:
    \begin{itemize}
        \item subconjunto do dataset final
        \item limitações da configuração atual
    \end{itemize}
\end{itemize}

\section{Resultados Iniciais de Desempenho}
\label{subsec:desempenho_inicial}

\textbf{Ideias a abordar:}
\begin{itemize}
    \item Métricas iniciais:
    \begin{itemize}
        \item acurácia geral
    \end{itemize}
    \item Comparações:
    \begin{itemize}
        \item entre diferentes prompts
        \item entre modelos (se aplicável)
    \end{itemize}
    \item Observações:
    \begin{itemize}
        \item desempenho global
        \item variações relevantes
    \end{itemize}
    \item Inclusão de:
    \begin{itemize}
        \item tabela de resultados
        \item gráfico simples (opcional, mas recomendado)
    \end{itemize}
\end{itemize}

\section{Análise por Tipo de Questão}
\label{subsec:tipo_questao}

\textbf{Ideias a abordar:}
\begin{itemize}
    \item Classificação das questões:
    \begin{itemize}
        \item objetivas simples
        \item situacionais
        \item múltiplas respostas (se houver)
    \end{itemize}
    \item Comparação de desempenho entre os tipos
    \item Identificação de padrões:
    \begin{itemize}
        \item melhor desempenho em questões factuais
        \item maior dificuldade em cenários complexos
    \end{itemize}
    \item Relação com os conceitos da fundamentação
\end{itemize}

\section{Análise Qualitativa dos Erros}
\label{subsec:analise_erros}

\textbf{Ideias a abordar:}
\begin{itemize}
    \item Identificação de padrões de erro:
    \begin{itemize}
        \item escolha de alternativa plausível, porém incorreta
        \item interpretação superficial do cenário
        \item respostas genéricas
    \end{itemize}
    \item Apresentação de exemplos:
    \begin{itemize}
        \item 1 ou 2 questões analisadas em detalhe
    \end{itemize}
    \item Classificação dos erros:
    \begin{itemize}
        \item erro factual
        \item erro de interpretação
        \item erro de julgamento
    \end{itemize}
\end{itemize}

\section{Efeito das Estratégias de Prompting}
\label{subsec:prompt_resultados}

\textbf{Ideias a abordar:}
\begin{itemize}
    \item Comparação entre estratégias:
    \begin{itemize}
        \item zero-shot vs outras abordagens
    \end{itemize}
    \item Impacto observado:
    \begin{itemize}
        \item melhoria (ou não) na acurácia
        \item impacto na consistência das respostas
    \end{itemize}
    \item Observações preliminares:
    \begin{itemize}
        \item dependência do tipo de questão
    \end{itemize}
\end{itemize}

\section{Discussão dos Resultados Preliminares}
\label{subsec:discussao_preliminar}

\textbf{Ideias a abordar:}
\begin{itemize}
    \item Interpretação geral dos resultados:
    \begin{itemize}
        \item coerência com o comportamento esperado dos LLMs
        \item alinhamento com os conceitos da fundamentação
    \end{itemize}
    \item Principais insights:
    \begin{itemize}
        \item limitações em cenários situacionais
        \item influência do prompt
    \end{itemize}
    \item Implicações iniciais:
    \begin{itemize}
        \item uso cauteloso de LLMs como apoio ao estudo para PMP
    \end{itemize}
\end{itemize}

\section{Limitações dos Resultados Preliminares}
\label{subsec:limitacoes_prelim}

\textbf{Ideias a abordar:}
\begin{itemize}
    \item Tamanho reduzido do dataset
    \item Cobertura limitada de domínios PMP
    \item Número reduzido de execuções
    \item Simplificação das estratégias de prompting
\end{itemize}

\section{Considerações Finais do Capítulo}
\label{subsec:resultados_final}

\textbf{Ideias a abordar:}
\begin{itemize}
    \item Síntese dos resultados preliminares
    \item Evidência de viabilidade do estudo
    \item Reforço da relevância da investigação
    \item Conexão com o próximo capítulo (ex.: conclusões ou trabalhos futuros)
\end{itemize}